\documentclass{article}
\usepackage[margin=2.5cm]{geometry}
\usepackage{amsmath,amssymb}
\usepackage{booktabs}

\title{Resource-Constrained Project Scheduling Problem \\ with Alternative Subgraphs (RCPSP-AS)}
\author{}
\date{}

\begin{document}

\maketitle

\section{Problem Definition}

A project is represented as a directed acyclic graph $G = (N, A)$ where $N = \{0, \ldots, n\}$ is a set of non-preemptive, topologically ordered activities and $A \subset N \times N$ is a set of end--start precedence arcs. Each activity~$i$ has a fixed duration~$d_i$ and demands $r_{i,v}$ units of each renewable resource $v \in R$, which has limited capacity~$a_v$. Activity~$0$ is the unique source and activity~$n$ the unique sink.

RCPSP-AS extends RCPSP with a set of \emph{alternative subgraphs}~$L$. Each subgraph $l \in L$ has a principal activity~$p_l$ (source of the subgraph), a terminal activity~$t_l$ (sink), and a set of branches~$K_{p_l}$. Each branch $k \in K_{p_l}$ starts from a \emph{branching activity}~$b_k$, an immediate successor of~$p_l$, and contains activities $N_{b_k} \subset N$ with $b_k \in N_{b_k}$, $p_l \notin N_{b_k}$, and $t_l \notin N_{b_k}$. Activities belonging to at least one branch form the \emph{alternative} set $N^a = \bigcup_k N_{b_k}$; the rest are \emph{fixed}: $N^f = N \setminus N^a$. Exactly one branch must be selected from each subgraph whose principal activity is scheduled.

Two extensions are supported: \emph{nested subgraphs} (a subgraph placed inside another's branch) and \emph{linked branches}, indicated by $\kappa_{i,j} = 1$ when activity $i \in N_{b_k}$ links to $j \in N_{b_{k'}}$ ($k \neq k'$), causing the selection of~$i$ to imply the selection of~$j$ and its successors in branch~$k'$.

The objective is to minimize the makespan $C_{\max} = s_n + d_n$.

\section{Notation}

\begin{table}[h]
\centering
\small
\begin{tabular}{@{}cl@{}}
\toprule
\textbf{Symbol} & \textbf{Description} \\
\midrule
\multicolumn{2}{@{}l}{\textit{Sets and indices}} \\[2pt]
$N = \{0, \ldots, n\}$ & Activities (topologically ordered); $0$ = source, $n$ = sink \\
$A \subset N \times N$ & Precedence arcs \\
$A_{\mathrm{prop}} = A \setminus \{(p_l, j) : l \in L\}$ & Non-branching arcs (simplified formulation) \\
$L$ & Alternative subgraphs \\
$K_{p_l}$ & Branches of subgraph $l$ \\
$N_{b_k} \subset N$ & Activities in branch $k$ \\
$N^a = \bigcup_k N_{b_k}$ & Alternative activities \\
$N^f = N \setminus N^a$ & Fixed activities \\
$R$ & Renewable resources \\
$\mathcal{T} = \{0, \ldots, T_{\max}\}$ & Planning horizon (ILP) \\
$K = \{k^*\} \cup \bigcup_{l} K_{p_l}$ & All branches including dummy $k^*$ (original CP) \\
$B_i = \{k \in K \mid i \in N_{b_k}\}$ & Branches containing activity $i$ (original CP) \\[2pt]
\midrule
\multicolumn{2}{@{}l}{\textit{Parameters}} \\[2pt]
$d_i$ & Duration of activity $i$ \\
$r_{i,v}$ & Demand of activity $i$ for resource $v$ \\
$a_v$ & Capacity of resource $v$ \\
$p_l$ & Principal activity of subgraph $l$ \\
$t_l$ & Terminal activity of subgraph $l$ \\
$b_k$ & Branching activity of branch $k$ \\
$\kappa_{i,j} \in \{0,1\}$ & Link indicator: 1 if $i$ in branch $k$ links to $j$ in branch $k' \neq k$ \\
$M$ & Sufficiently large constant (ILP) \\[2pt]
\midrule
\multicolumn{2}{@{}l}{\textit{Decision variables --- ILP}} \\[2pt]
$x_i \in \{0, 1\}$ & 1 if activity $i$ is selected \\
$s_i \in \mathbb{Z}_{\ge 0}$ & Start time of activity $i$ \\[2pt]
\midrule
\multicolumn{2}{@{}l}{\textit{Decision variables --- CP}} \\[2pt]
$\mathsf{x}_i$ & Optional interval variable for activity $i$, size $d_i$ \\
$\mathrm{presenceOf}(\mathsf{x}_i)$ & 1 if interval $\mathsf{x}_i$ is present, 0 otherwise \\
$\mathrm{startOf}(\mathsf{x}_i)$ & Start time of interval $\mathsf{x}_i$ \\
$\mathrm{endOf}(\mathsf{x}_i)$ & End time of interval $\mathsf{x}_i$ \\
$\mathrm{pulse}(\mathsf{x}_i, h)$ & Contributes $h$ to a cumulative function while $\mathsf{x}_i$ is present \\
\bottomrule
\end{tabular}
\end{table}


%% ============================================================================
\section{Original Formulation}
%% ============================================================================

This formulation, following Servranckx \& Vanhoucke (2019), distinguishes fixed and alternative activities explicitly and uses dedicated constraints for branch successor implications and inter-branch links.

\subsection{ILP}

\begin{align}
\min \quad & s_n + d_n \label{eq:o-obj} \\[2mm]
\text{s.t.} \quad
& x_i = 1
    && \forall i \in N^f \label{eq:o-fixed} \\
& \sum_{k \in K_{p_l}} x_{b_k} = x_{p_l}
    && \forall l \in L \label{eq:o-branch} \\
& x_{p_l} \le x_{t_l}
    && \forall l \in L \label{eq:o-terminal} \\
& x_i \le x_j
    && \forall l \in L,\; k \in K_{p_l},\; (i,j) \in (N_{b_k} \!\times\! N_{b_k}) \cap A \label{eq:o-succ} \\
& x_i \le x_j
    && \forall (i,j) :\, \kappa_{i,j} = 1 \label{eq:o-link} \\
& s_i + d_i \le s_j + M(2 - x_i - x_j)
    && \forall (i,j) \in A \label{eq:o-prec} \\
& \sum_{\substack{i \in N:\\ s_i \le t < s_i\!+\!d_i}} r_{i,v} \cdot x_i \le a_v
    && \forall v \in R,\; t \in \mathcal{T} \label{eq:o-res} \\
& x_i \in \{0,1\},\;\; s_i \in \mathbb{Z}_{\ge 0}
    && \forall i \in N \label{eq:o-var}
\end{align}

\noindent
\textbf{\eqref{eq:o-fixed}} All fixed activities are always selected.
\textbf{\eqref{eq:o-branch}} Exactly one branch is chosen per active subgraph.
\textbf{\eqref{eq:o-terminal}} If the principal activity is selected, so is the terminal.
\textbf{\eqref{eq:o-succ}} Within a branch, selection propagates along precedence arcs.
\textbf{\eqref{eq:o-link}} Linked activities propagate selection across branches.
\textbf{\eqref{eq:o-prec}} Big-$M$ linearization of precedence: if both $i,j$ are selected, $i$ finishes before $j$ starts.
\textbf{\eqref{eq:o-res}} Resource capacity is respected at every time point.


\subsection{CP}

Before building the model, a preprocessing step eliminates linked branches and unifies activity selection:
%
\begin{enumerate}
    \item \textbf{Eliminate links.} For each $\kappa_{i,j} = 1$ between branches $k$ and $k'$, extend $N_{b_k} \leftarrow N_{b_k} \cup \bigl(S^*_j \cap N_{b_{k'}}\bigr)$, where $S^*_j$ denotes the transitive successors of~$j$.
    \item \textbf{Dummy branch.} Introduce $k^*$ with $b_{k^*} = 0$ and $N_{b_{k^*}} = N^f$.
    \item \textbf{Derived sets.} $K = \{k^*\} \cup \bigcup_{l \in L} K_{p_l}$ \;\; and \;\; $B_i = \{k \in K \mid i \in N_{b_k}\}$.
\end{enumerate}

\begin{align}
\min \quad & \mathrm{endOf}(\mathsf{x}_n) \label{eq:ocp-obj} \\[2mm]
\text{s.t.} \quad
& \mathrm{presenceOf}(\mathsf{x}_0) = 1
    \label{eq:ocp-source} \\
& \mathrm{presenceOf}(\mathsf{x}_i) \;\Leftrightarrow\; \bigvee_{k \in B_i} \mathrm{presenceOf}(\mathsf{x}_{b_k})
    && \forall i \in N \!\setminus\! \{0\} \label{eq:ocp-sel} \\
& \sum_{k \in K_{p_l}} \mathrm{presenceOf}(\mathsf{x}_{b_k}) = \mathrm{presenceOf}(\mathsf{x}_{p_l})
    && \forall l \in L \label{eq:ocp-branch} \\
& \bigl(\mathrm{presenceOf}(\mathsf{x}_i) \!\wedge\! \mathrm{presenceOf}(\mathsf{x}_j)\bigr) \Rightarrow \mathrm{endOf}(\mathsf{x}_i) \le \mathrm{startOf}(\mathsf{x}_j)
    && \forall (i,j) \!\in\! A \label{eq:ocp-prec} \\
& \sum_{i \in N} \mathrm{pulse}(\mathsf{x}_i,\, r_{i,v}) \le a_v
    && \forall v \in R \label{eq:ocp-res} \\
& \mathsf{x}_i \;\text{optional interval, size}\; d_i
    && \forall i \in N \label{eq:ocp-var}
\end{align}

\noindent
After preprocessing, activity selection~\eqref{eq:ocp-sel} is unified: activity~$i$ is present if and only if at least one branch containing it is selected. This single rule covers both fixed activities (via the dummy branch~$k^*$ tied to the always-present source) and alternative activities. Constraints \eqref{eq:ocp-branch}--\eqref{eq:ocp-res} mirror the ILP: branch selection, conditional precedence timing, and cumulative resource limits.

%% ============================================================================
\section{Simplified Formulation}
%% ============================================================================

The simplified formulation replaces the explicit classification into $N^f$, $N^a$, $t_l$, and $\kappa$ with a single mechanism: \emph{selection propagation} along the set of non-branching arcs $A_{\mathrm{prop}} = A \setminus \{(p_l, j) \in A : l \in L\}$. Since the source is always selected and $A_{\mathrm{prop}}$ contains all arcs except those from principal activities to their branch successors, fixed activities are selected by transitivity. Within a selected branch, successor selection follows the same propagation. Cross-branch links are handled automatically because the link arc $(i, j)$ belongs to $A_{\mathrm{prop}}$ whenever the source~$i$ is not a principal activity. No preprocessing is needed.

\subsection{ILP}

\begin{align}
\min \quad & s_n + d_n \label{eq:s-obj} \\[2mm]
\text{s.t.} \quad
& x_0 = 1
    \label{eq:s-source} \\
& x_i \le x_j
    && \forall (i,j) \in A_{\mathrm{prop}} \label{eq:s-prop} \\
& \sum_{j:\,(p_l,j) \in A} x_j = x_{p_l}
    && \forall l \in L \label{eq:s-branch} \\
& s_i + d_i \le s_j + M(2 - x_i - x_j)
    && \forall (i,j) \in A \label{eq:s-prec} \\
& \sum_{\substack{i \in N:\\ s_i \le t < s_i\!+\!d_i}} r_{i,v} \cdot x_i \le a_v
    && \forall v \in R,\; t \in \mathcal{T} \label{eq:s-res} \\
& x_i \in \{0,1\},\;\; s_i \in \mathbb{Z}_{\ge 0}
    && \forall i \in N \label{eq:s-var}
\end{align}

\noindent
\textbf{\eqref{eq:s-source}} The source activity is selected.
\textbf{\eqref{eq:s-prop}} Selection propagates along non-branching arcs---this single rule replaces~\eqref{eq:o-fixed}, \eqref{eq:o-terminal}, \eqref{eq:o-succ}, and~\eqref{eq:o-link}.
\textbf{\eqref{eq:s-branch}} Exactly one successor of $p_l$ is chosen iff $p_l$ is selected.
\textbf{\eqref{eq:s-prec}}--\textbf{\eqref{eq:s-res}} Precedence and resource constraints are unchanged.


\subsection{CP}

\begin{align}
\min \quad & \mathrm{endOf}(\mathsf{x}_n) \label{eq:scp-obj} \\[2mm]
\text{s.t.} \quad
& \mathrm{presenceOf}(\mathsf{x}_0) = 1
    \label{eq:scp-source} \\
& \mathrm{presenceOf}(\mathsf{x}_i) \Rightarrow \mathrm{presenceOf}(\mathsf{x}_j)
    && \forall (i,j) \in A_{\mathrm{prop}} \label{eq:scp-prop} \\
& \sum_{j:\,(p_l,j) \in A} \mathrm{presenceOf}(\mathsf{x}_j) = \mathrm{presenceOf}(\mathsf{x}_{p_l})
    && \forall l \in L \label{eq:scp-branch} \\
& \bigl(\mathrm{presenceOf}(\mathsf{x}_i) \!\wedge\! \mathrm{presenceOf}(\mathsf{x}_j)\bigr) \Rightarrow \mathrm{endOf}(\mathsf{x}_i) \le \mathrm{startOf}(\mathsf{x}_j)
    && \forall (i,j) \!\in\! A \label{eq:scp-prec} \\
& \sum_{i \in N} \mathrm{pulse}(\mathsf{x}_i,\, r_{i,v}) \le a_v
    && \forall v \in R \label{eq:scp-res} \\
& \mathsf{x}_i \;\text{optional interval, size}\; d_i
    && \forall i \in N \label{eq:scp-var}
\end{align}

\noindent
The CP formulation mirrors the simplified ILP directly. Optional interval variables~\eqref{eq:scp-var} unify the selection ($\mathrm{presenceOf}$) and timing ($\mathrm{startOf}$, $\mathrm{endOf}$) decisions. The $\mathrm{pulse}$ function~\eqref{eq:scp-res} models time-varying resource usage, enabling the solver's built-in cumulative constraint.


\section{References}

\noindent
Servranckx, T., \& Vanhoucke, M. (2019). A tabu search procedure for the resource-constrained project scheduling problem with alternative subgraphs. \textit{European Journal of Operational Research}, 273(3), 841--860.

\end{document}
